\chapter{Data Generating Function}

\begin{definition}[Data generating function]
	\textcolor{red}{($\star$) TBD}
\end{definition}

If we have a random variable $T$ with parameter $\theta$, a data generating function (DGF) with respect to $T$ is a function $\tau$, 
\(
    \tau(\mathbf{U}, \theta) 
\)
such that $T \overset{d}{=} \tau(\mathbf{U}, \theta) $. $\mathbf{U}$ is a random vector independent of the parameter $\theta$. Intuitively, one can say that
\[
\text{data} = \text{randomness} + \text{parameter}
\]
and the parameter gets abgetrennt from the randomness. The random vector $\mathbf{U}$ can not depend on $\theta$. So randomness and the parameter $\theta$ are independent. 

\[
    \mathbf{U} \overset{\tau(., \theta)}{ \longrightarrow} T
\]

Example:
\begin{enumerate}
    \item Normal distribution
    \item wo es keine DGF gibt
\end{enumerate}


\begin{example}[Example: DGF, Normal distribution] \label{example: DGF normal} 
	Let $T \sim \mathcal{N}(\theta, 1)$. Then we can say that 
	\[
	T \sim \theta + U
	\]
	where $U \sim \mathcal{N}(0, 1)$. So $U$ does not depend on the parameter $\theta$. We can define the DGF
	\begin{equation*}
		\tau(U, \theta) = \theta + U.
	\end{equation*}
	
	If $\theta$ is given, the DGF $\tau$ generates simulated data of $T$. Assume we observe $T= t$ and $U = u$ is given. Then we can get a best estimate of the parameter $\theta$ by inverting $\tau$:
	\begin{align*}
		t &= \tau(u, \theta) = \theta + u \\
		\implies \theta &= t - u
	\end{align*}
	So the inversion is $\hat \theta (u, t) = t-u$. 
	
	In this example, the inversion works perfectly, because $\tau$ is as a function of $\theta$ bijective and therefore invertible. There are also DGFs which can not be inverted, as one can see in Example \ref{example: non invertible DGF}.
	
\end{example}

\begin{example}[Example: DGF with $n > 1$]
	We can generalize Example \ref{example: DGF normal} by defining
	$T = \sum_{i=1}^n X_i$ where $X_1, \ldots , X_n \overset{iid}{\sim} \mathcal{N}(\theta, 1)$.  Again use the interpretation as in Example \ref{example: DGF normal}: $X_i = \theta + U_i$ for $i \in \{1, \ldots, n\}$ with $U_1, \ldots , U_n \overset{iid}{\sim} \mathcal{N}(0, 1)$. Then it holds $T = \sum_{i=1} ^n \theta + U_i = n \theta + \sum_{i=1} ^n U_i$, so the DGF is defined as 
	\[
	    \tau (\mathbf{U}, \theta )  = n \theta + \sum_{i=1} ^n U_i.
	\]
	This again is invertible with
	\[
	    \hat \theta(\mathbf{u}, t ) = \frac{t - \sum_{i=1} ^n u_i}{n}.
	\]
	
	Figure \textcolor{red}{($\star$) insert Figure} shows the DGF $\tau(\mathbf{U}, .)$.
	
	It is a strictly monotone function \textcolor{red}{($\star$) Is this function really strictly monotone or something else?) and therefore bijective, so we can always find a $\theta$ given $T$ and $U$.}
\end{example}

\begin{example}[Non invertible DGF] \label{example: non invertible DGF}
	
	Let $\theta \in [0, 1]$, $T \sim Bin(n , \theta)$ and $U_1, \ldots, U_n \overset{iid}{\sim} U[0, 1]$. It holds 
	\[
	    T \sim \sum_{i=1}^n   \mathds{1} ( U_i < \theta )
	\]
	
	We can see this by finding the CDF of $X_i := \mathds{1} ( U_i < \theta )$. 
	$X_i$ can take values of 0 and 1. $X_i$ takes the value 1 if and only if $U_i < \theta$, so
	
	\begin{align*}
	    \mathbb P (X_i = 1) 
	    = 
	    \mathbb P (U_i < \theta)  = \theta
	\end{align*}
	where the last equality comes from the fact that $U_i$ is standard uniformly distributed. 
	So it holds $X_i \sim Bin(1, \theta)$ and since the sum is again binomial, it holds
	\[
	    \sum_{i = 1}^n \underbrace{\mathds{1} ( U_i < \theta )}_{\sim Bin(1, \theta)} \sim Bin(n, \theta)
	\]
	and therefore
	\[
	    T \sim \sum_{i = 1}^n \mathds{1} ( U_i < \theta ).
	\]
	
	So we can define the DGF
	\[
	    \tau (\mathbf{U}, \theta) = \sum_{i=1}^n \mathds{1} ( U_i < \theta ).
	\]
	
	Fix $n = 2, \mathbf{U} = (0.2, 0.8), T = 1$. 
	Plotting $\tau(\mathbf{U}, .)$ gives Figure \textcolor{red}{($\star$) Figure einfügen}
	
	The equation
	\begin{align*}
	      t &= \mathds{1} (U_1 < \theta) + \mathds{1} (U_2 < \theta)& \\
	    \implies
	    1 &= \mathds{1} (0.2 < \theta) + \mathds{1} (0.8 < \theta)&
	\end{align*}
	  
	
	delivers solutions for $\theta$: $\theta \in (0.2, 0.8]$. So there are infinite many solutions. In particular the solution is not unique. So the DGF $\tau$ in not invertible.

\end{example}


Theorem 1 in \cite{Lindqvist2025simulated} is given by 
\begin{theorem}\label{theorem:dgf-upper bound}
	Define \[\bar \theta ( \mathbf u, t) = \sup \{ \theta: \tau (\mathbf u, \theta) \leq t \}\}\] and let $a(t)$ be a left continuous function such that
	\[
		\mathbb{P} \left( \bar \theta (\mathbf U, t)  \geq a(t)\right) \leq \alpha
	\]
	for all $t$ and $\alpha \in (0, 1)$. 
	Then $a(T)$ is an upper confidence bound for $\theta$ that is $\mathbb{P} \left( \theta \geq a(T) \right) \leq \alpha $ for all $\alpha \in (0, 1)$.
\end{theorem}

\begin{example}
	Suppose $U \sim \text{Exp}(1)$ and $T \sim \text{Exp} \left( \frac{1}{\theta} \right)$. Then we can express $T$ in terms of the function $\tau$ if we define it as 
	\[ \tau (u , \theta ) = u \cdot \theta\].
	because then 
	\[
		\mathbb{P} (\tau (U , \theta ) \leq x) = 
		\mathbb{P} \left( U \cdot \theta \leq x \right) =
		\mathbb{P} \left( U \leq \frac{x}{\theta} \right) = 
		1 - \exp \left( - \frac{x}{\theta} \right)
	\]	
	and this is the CDF of an exponentially distributed random variable with parameter $ \frac{x}{\theta}$. So it holds $T \overset{d}{=} \tau (U, \theta)$. Now we caculate $\bar \theta (u, t)$:
	\[
		\bar \theta (u, t) = 
		\sup \left\{ \theta : \tau (u, \theta) \leq t \right\}
		=
		\sup \left\{ \theta : u \cdot \theta \leq t \right\}
		=
		\sup \left\{ \theta : \theta \leq \frac{t}{u} \right\}
		=
		\frac{t}{u}.
	\]
	Let $\alpha \in (0, 1)$ be arbitrary.
	Now the next objective is to find the function $a$ such that $\mathbb{P} \left( \bar \theta (u, t) \geq a(t) \right) = \alpha$.
	\begin{align*}
		&\mathbb{P} \left( \bar \theta (u, t) \geq a(t) \right) = \alpha
		\\
		\Leftrightarrow &
		\mathbb{P} \left( \frac{t}{u} \geq a(t) \right) = \alpha 
		\\
		\Leftrightarrow &
		\mathbb{P} \left( u \leq \frac{t}{a(t)} \right) = \alpha 
		\\
		\Leftrightarrow &
		1 - \exp \left( - \frac{t}{a(t)} \right) = \alpha 
		\\
		\Leftrightarrow &
		a(t) = - \frac{t}{\ln(1- \alpha)}
	\end{align*}
	This function is left continuous because \textcolor{red}{($\star$) tbd} so therefore we can use it for Theorem~\ref{theorem:dgf-upper bound}.
	So after observing $T$, one can say that 
	\[
		\left( 0, a(T) \right] = \left( 0,  \frac{-T}{\ln(1- \alpha)} \right]
	\]
	is a one-sided confidence interval for $\theta$. For example for $\alpha = 0.1$ and the upper confidence bound is given by $9.49 T$.
\end{example}

\begin{proof}
	Let us define 
	\[
		a^{-1} (\theta ) = \sup \{ t : a(t) \leq \theta \}.
	\]
	First, observe that if $\theta \geq a(T)$, then it also holds $T \leq a^{-1}(T)$ and consequently $\{ \theta \geq a(T) \} \subseteq \{ T \leq a^{-1}(T) \}$. To see this, assume $\theta \geq a(T)$. Then $T \in \{ t: a(t) \leq \theta \}$. Since $T$ is an element of this set, $T$ has to be smaller or equal to the supremum of this set, i.e. 
	\[
		T \leq \sup \{ t: a(t) \leq \theta \}
	\]
	This is a pure consequence of the definition of the supremum of a set. The supremum is exactly $a^{-1}(\theta)$ and therefore it holds $T \leq a^{-1}(\theta)$. All in all, the inequality $\theta \geq a(T)$, implies $T \leq a^{-1}(T)$. Since the set  $\{ \theta \geq a(T) \}$ is just a subset of $\{ T \leq a^{-1}(T) \}$, it holds
	\[
		\mathbb{P} \left( \theta \geq a(T) \right) \leq 
		\mathbb{P} \left( T \leq a^{-1} (\theta) \right).
	\]
	Next, we use the definition of $\tau$
	\[
		\mathbb{P} \left( T \leq a^{-1} (\theta) \right)
		=
		\mathbb{P} \left( \tau(\mathbf U, \theta) \leq a^{-1} (\theta) \right).
	\]
	Next, we show that $\{ \tau ( \mathbf U, \theta ) \leq a^{-1}(\theta) \} \subseteq \{ \bar \theta (\mathbf U, a^{-1}(\theta)) \geq \theta \}$. For this assume $ \tau ( \mathbf U, \theta ) \leq a^{-1}(\theta)$. Define $t' := a^{-1}(\theta)$, then it holds $ \tau ( \mathbf U, \theta ) \leq t'$. So, we can observe that $\theta \in \{\theta' : \tau(\mathbf U, \theta') \leq t'\}$. Again, with the properties of the supremum, we get
	\[
		\theta \leq \sup \{\theta' : \tau(\mathbf U, \theta') \leq t'\}
	\]
	and the right hand side of this inequality is exactly $\bar \theta (\mathbf U, t')$. All in all, it holds that $\tau (\mathbf U, \theta) \leq t'$ implies that $\theta \leq \bar \theta (\mathbf U, t')$. We can use this and obtain
	\[
		\mathbb{P} \left( \tau (\mathbf U, \theta) \leq a^{-1} (\theta) \right)
		\leq 
		\mathbb{P} \left( \bar \theta \left( \mathbf U, a^{-1}(\theta) \right)  \geq \theta \right).
	\]
	
	\paragraph{Step (iv): It holds $a\big(a^{-1}(\theta)\big) \le \theta$, hence
		$\{\theta(U,a^{-1}(\theta)) \ge \theta\} \subseteq \{\theta(U,a^{-1}(\theta)) \ge a(a^{-1}(\theta))\}$.}
	
	This is where left-continuity of $a$ is used. Let
	\[
	t^\ast := a^{-1}(\theta) = \sup\{t : a(t) \le \theta\}.
	\]
	
	\begin{itemize}
		\item If the defining set $\{t : a(t) \le \theta\}$ has a maximum, i.e.\ contains $t^\ast$ itself, then trivially $a(t^\ast) \le \theta$ and we are done.
		\item Otherwise, $t^\ast$ is only a limit point of the set: there is a sequence $t_n \uparrow t^\ast$ with $a(t_n) \le \theta$ for every $n$. Left-continuity of $a$ means
		\[
		a(t_n) \longrightarrow a(t^\ast) \quad \text{as } t_n \uparrow t^\ast.
		\]
		Since $a(t_n) \le \theta$ for all $n$, and the limit of a sequence bounded above by $\theta$ is itself $\le \theta$, we obtain
		\[
		a(t^\ast) \le \theta.
		\]
	\end{itemize}
	
	So
	\[
	a\big(a^{-1}(\theta)\big) \le \theta.
	\]
	
	Now let $X := \bar \theta\big(U, a^{-1}(\theta)\big)$. The threshold on the right-hand side, $a(a^{-1}(\theta))$, is no larger than the threshold $\theta$ on the left-hand side. A weaker (smaller) threshold makes the event ``$X \ge \text{threshold}$'' easier to satisfy, i.e.\ the event is larger:
	\[
	\{X \ge \theta\} \subseteq \big\{X \ge a(a^{-1}(\theta))\big\},
	\]
	since $a(a^{-1}(\theta)) \le \theta \le X$ whenever $X \ge \theta$. Taking probabilities on both sides gives
	\[
	\mathbb{P}_U\Big( \bar\theta\big(U, a^{-1}(\theta)\big) \ge \theta \Big)
	\;\le\;
	\mathbb{P}_U\Big( \bar\theta\big(U, a^{-1}(\theta)\big) \ge a\big(a^{-1}(\theta)\big) \Big).
	\]
	
	
	\paragraph{Step (v): $\mathbb{P}_U\big(\bar\theta(U,a^{-1}(\theta)) \ge a(a^{-1}(\theta))\big) \le \alpha$.}
	
	This is now a direct instance of the hypothesis
	\[
	\mathbb{P}_U\big( \theta(U,t) \ge a(t) \big) \le \alpha \qquad \text{for all } t,
	\]
	applied with the particular choice $t = a^{-1}(\theta)$ --- a legitimate substitution since the hypothesis is assumed to hold \emph{for all} $t \in \mathbb{R}$, and $a^{-1}(\theta)$ is simply one specific real number, deterministic given $\theta$. Hence
	\[
	\mathbb{P}_U\Big( \theta\big(U, a^{-1}(\theta)\big) \ge a\big(a^{-1}(\theta)\big) \Big) \le \alpha.
	\]
	
	Combining everything gives \textcolor{red}{($\star$) Alle ungleichungen zusammenpacken, eventuell die steps eindeutig kennzeichnen.}
\end{proof}



\section{Clopper Pearson Intervals}
\textcolor{red}{($\star$) hier example einführen und klar machen wie man die theorems und corollaries benutzt. Vielleicht ein Bild einfügen (von den ordnungsstatistiken U oder den graphen von $\tau$).}



